%=====================================================================
%  PSO_GA_explained.tex
%  A beginner-friendly explanation of Particle Swarm Optimization (PSO)
%  and the Genetic Algorithm (GA) IN THE CONTEXT OF THIS PROJECT:
%  "AI-Powered Autonomous Smart Agriculture using Swarm Intelligence
%   and Reinforcement Learning".
%
%  Compile with:  pdflatex PSO_GA_explained.tex
%=====================================================================
\documentclass[11pt]{article}

\usepackage[a4paper,margin=1in]{geometry}
\usepackage{amsmath}
\usepackage{amssymb}
\usepackage{booktabs}
\usepackage{xcolor}
\usepackage{tcolorbox}
\usepackage{enumitem}
\usepackage{hyperref}
\hypersetup{colorlinks=true, linkcolor=blue!60!black, urlcolor=blue!60!black}

% A soft coloured box for "intuition / analogy" notes.
\newtcolorbox{intuition}[1][]{colback=blue!4!white, colframe=blue!45!black,
  fonttitle=\bfseries, title=Intuition, #1}
\newtcolorbox{codebox}[1][]{colback=gray!6!white, colframe=gray!55!black,
  fonttitle=\bfseries, title=In our code, #1}
\newtcolorbox{whybox}[1][]{colback=green!4!white, colframe=green!40!black,
  fonttitle=\bfseries, title=Why we did this, #1}

\newcommand{\tab}{\hspace{2em}}

\title{\textbf{Population-Based Optimization for Autonomous Pollination}\\[4pt]
\large A Plain-Language Guide to PSO and the Genetic Algorithm\\ in the
Smart-Agriculture Project}
\author{AI Laboratory Project}
\date{}

\begin{document}
\maketitle

\begin{abstract}
\noindent
This document explains, in simple language, the two \emph{population-based
optimization} algorithms used in Phase~1 of the project: \textbf{Particle Swarm
Optimization (PSO)} and the \textbf{Genetic Algorithm (GA)}. Both algorithms
solve the \emph{same} question --- \textbf{which pollination robot should visit
which flowers?} --- but in two very different ways. PSO is inspired by a flock of
birds; GA is inspired by biological evolution. We describe what problem we are
solving, the exact idea behind each algorithm, every formula/operator we use, and
\emph{where and why} each design choice appears in our code. No prior knowledge
of optimization is assumed.
\end{abstract}

\tableofcontents

%=====================================================================
\section{The Problem We Are Solving (Phase 1: Task Allocation)}
%=====================================================================

Imagine a large greenhouse. Inside it there are:
\begin{itemize}[nosep]
  \item \textbf{400 flowers} that must all be pollinated,
  \item \textbf{5 pollination robots} that fly around and pollinate,
  \item a few \textbf{charging stations} and some \textbf{obstacles}.
\end{itemize}

Our job in Phase~1 is \textbf{task allocation}: decide, for every one of the
400 flowers, \emph{which} of the 5 robots is responsible for it. We are
\textbf{not} deciding \emph{how} a robot flies (that is Phase~2, done later by
Value Iteration and Q-Learning). We are only deciding the \emph{teams}: robot~0
gets these flowers, robot~1 gets those flowers, and so on.

\begin{intuition}
Think of a teacher splitting 400 exam papers among 5 teaching assistants. You
want the split to be \emph{fair} (nobody gets 300 papers while another gets 20)
and \emph{smart} (give each assistant the papers that are closest / easiest for
them). Task allocation is exactly this kind of ``divide the work well'' problem.
\end{intuition}

\subsection{What makes an allocation ``good''? --- the fitness function}

A good allocation should:
\begin{enumerate}[nosep]
  \item \textbf{Cover everything} --- all 400 flowers get pollinated.
  \item \textbf{Short travel} --- robots should not fly further than necessary.
  \item \textbf{Low energy} --- less flying means less battery used.
  \item \textbf{Finish fast} --- the busiest robot should still finish quickly.
  \item \textbf{Balanced workload} --- every robot does a similar share.
  \item \textbf{Use all robots} --- do not leave a robot sitting idle.
\end{enumerate}

We combine these six goals into a single number called the
\textbf{fitness} (or \emph{cost}). Each goal is written as a \emph{penalty}
that is $0$ when perfect and grows when things get worse:

\begin{align}
\text{fitness} =\; & 5.0\,(1-\text{coverage})
+ 1.0\,\frac{\text{distance}}{\text{ref}_d}
+ 1.0\,\frac{\text{energy}}{\text{ref}_e} \notag\\
& + 1.0\,\frac{\text{time}}{\text{ref}_t}
+ 1.0\,\frac{\text{workload\_std}}{\text{ref}_b}
+ 0.5\,(1-\text{utilisation}).
\label{eq:fitness}
\end{align}

\begin{tcolorbox}[colback=yellow!8!white,colframe=orange!60!black,title=Very important: lower fitness = better]
Because every term is a \textbf{penalty}, a \emph{smaller} fitness value means a
\emph{better} allocation. A perfect allocation would give fitness near $0$.
This is a \textbf{minimization} problem. (Later, in Phase~2, Q-Learning and
Value Iteration \emph{maximize} a reward instead --- opposite convention, same
idea.) So when we say ``GA scored $1.43$ and PSO scored $1.59$'', GA is
\emph{better} because $1.43 < 1.59$.
\end{tcolorbox}

The weights ($5.0, 1.0, \dots$) say how much we care about each goal.
\textbf{Coverage} has the biggest weight ($5.0$) because pollinating every
flower is the most important thing. The reference scales
($\text{ref}_d,\text{ref}_e,\dots$) simply divide each raw quantity so that all
six terms are roughly the same size ($\approx 1$) and no single term dominates
just because of its units (metres vs.\ seconds vs.\ battery units).

\begin{codebox}
This is the function \texttt{evaluate(position, env)} in \texttt{PSO.py}. Both
PSO \emph{and} GA call this \emph{same} function to score their allocations, so
the comparison between them is perfectly fair.
\end{codebox}

%=====================================================================
\section{How Do We Describe an ``Allocation''?}
%=====================================================================

Before any algorithm can search for a good allocation, we need a way to
\emph{write down} an allocation as data. We use a list of 400 numbers:

\[
\text{allocation} = [\,r_0,\; r_1,\; r_2,\; \dots,\; r_{399}\,],
\qquad r_i \in \{0,1,2,3,4\}.
\]

Here $r_i$ is the robot assigned to flower $i$. For example
$[\,2,0,0,4,1,\dots]$ means ``flower~0 $\to$ robot~2, flower~1 $\to$ robot~0,
flower~2 $\to$ robot~0, \dots''. This single list of 400 robot-indices is one
complete candidate solution.

\begin{intuition}
One allocation = one complete ``answer'' to the whole problem. There are
$5^{400}$ possible answers (an astronomically huge number), so we cannot try
them all. PSO and GA are two clever ways to \emph{search} this huge space
without checking every possibility.
\end{intuition}

%=====================================================================
\section{Algorithm 1 --- Particle Swarm Optimization (PSO)}
%=====================================================================

\subsection{The core idea (bird flock analogy)}

PSO imitates a \textbf{flock of birds} searching for food. Each bird (called a
\textbf{particle}) flies around the search space. Every bird remembers the best
spot \emph{it personally} has ever found, and the whole flock shares the best
spot \emph{anyone} has found. Each bird keeps adjusting its flight to head partly
toward its own best memory and partly toward the flock's best. Over time the
whole flock converges on a very good spot.

\begin{intuition}
``Follow your own best experience, but also follow the leader.'' That single
sentence is the whole of PSO.
\end{intuition}

\subsection{What is a particle in \emph{our} problem?}

\begin{tcolorbox}[colback=blue!5!white,colframe=blue!50!black,title=One particle represents all 400 flowers]
A particle is a vector of \textbf{400 numbers} (one number per flower). But these
numbers are \emph{continuous} (real numbers like $2.73$), not robot indices yet.
Element $i$ is a real value in the range $[0,5)$. We convert it to a robot index
by taking the floor:
\[
\text{robot for flower } i = \big\lfloor \text{position}_i \big\rfloor,
\qquad\text{e.g. } \lfloor 2.73 \rfloor = 2.
\]
So a particle at position $[2.73,\,0.11,\,0.95,\,4.20,\dots]$ decodes to the
allocation $[2,0,0,4,\dots]$. This is called \textbf{continuous encoding of a
discrete problem}.
\end{tcolorbox}

\begin{codebox}
\texttt{decode(position, n\_robots)} does
$\lfloor \text{position} \rfloor$ then clips to $[0,4]$. The swarm is created
with \texttt{pos = rng.uniform(0, 5, size=(swarm\_size, n\_flowers))}, i.e.\
\texttt{swarm\_size} particles, each a 400-dimensional vector.
\end{codebox}

\subsection{How a particle moves --- the two PSO equations}

Each particle has a \textbf{position} $\mathbf{x}$ (where it currently is) and a
\textbf{velocity} $\mathbf{v}$ (which way and how fast it is moving). Every
iteration we update velocity, then position:

\begin{align}
\mathbf{v} &\leftarrow \underbrace{w\,\mathbf{v}}_{\text{inertia}}
 + \underbrace{c_1 r_1 (\mathbf{p}_{\text{best}} - \mathbf{x})}_{\text{cognitive: pull to own best}}
 + \underbrace{c_2 r_2 (\mathbf{g}_{\text{best}} - \mathbf{x})}_{\text{social: pull to flock best}}
 \label{eq:vel}\\[4pt]
\mathbf{x} &\leftarrow \mathbf{x} + \mathbf{v}
 \label{eq:pos}
\end{align}

Reading equation~\eqref{eq:vel} in plain words:
\begin{itemize}[nosep]
  \item $w\,\mathbf{v}$ \textbf{(inertia)}: keep some of your current momentum,
        so you do not stop or turn too abruptly. We use $w=0.7$.
  \item $c_1 r_1(\mathbf{p}_{\text{best}}-\mathbf{x})$ \textbf{(cognitive)}: get
        pulled back toward the best position \emph{you} have personally found.
        We use $c_1=1.5$.
  \item $c_2 r_2(\mathbf{g}_{\text{best}}-\mathbf{x})$ \textbf{(social)}: get
        pulled toward the best position \emph{the whole swarm} has found. We use
        $c_2=1.5$.
  \item $r_1,r_2$ are fresh random numbers in $[0,1]$ every step --- they add
        exploration so the swarm does not march in lockstep.
\end{itemize}

\begin{intuition}
$\mathbf{p}_{\text{best}}$ = ``the best place I remember.''\\
$\mathbf{g}_{\text{best}}$ = ``the best place anyone has told me about.''\\
A particle blends: keep going a bit (inertia) $+$ drift toward my memory $+$
drift toward the leader. Do this for every particle, every iteration, and the
flock closes in on a great allocation.
\end{intuition}

\begin{whybox}
We also \textbf{clamp velocity} (\texttt{vmax\_frac=0.2}) so particles do not fly
out of the useful region, \textbf{reflect} particles that leave the range
$[0,5)$ back inside, and \textbf{stop early} if the best fitness has not improved
for 30 iterations. These keep the search stable and efficient.
\end{whybox}

\begin{codebox}
This is \texttt{run\_pso(...)} in \texttt{PSO.py}. Hyperparameters:
\texttt{swarm\_size=40, inertia=0.7, c1=1.5, c2=1.5, max\_iter=120,
vmax\_frac=0.2, patience=30}. A fixed seed (\texttt{SEED=42}) makes every run
produce \emph{identical} results (reproducibility).
\end{codebox}

\subsection{The honest weakness of PSO here}
Our problem has \textbf{400 dimensions} (one per flower). Continuous PSO is
excellent in low dimensions but struggles when the number of dimensions is very
large, because neighbouring flowers are treated as independent numbers with no
built-in notion that ``nearby flowers should go to the same robot.'' As a
result, PSO's allocation is only modestly better than random. \emph{This weakness
is exactly why we also built a Genetic Algorithm allocator}, which uses a
discrete encoding and a heuristic starting point to reach shorter routes and
more balanced teams.

%=====================================================================
\section{Algorithm 2 --- Genetic Algorithm (GA)}
%=====================================================================

\subsection{The core idea (survival of the fittest)}

A Genetic Algorithm imitates \textbf{biological evolution}. It keeps a
\textbf{population} of candidate solutions. The \emph{fittest} solutions are more
likely to become ``parents.'' Parents are \textbf{combined} (crossover) to
create children, and children are \textbf{randomly tweaked} (mutation). Over many
\emph{generations}, the population evolves toward better and better solutions ---
just as useful traits spread through a species over time.

\begin{intuition}
``Keep the best answers, breed them together, occasionally mutate, and repeat.''
Good traits (short routes, balanced loads) survive and spread; bad ones die out.
\end{intuition}

\subsection{What is a chromosome in \emph{our} problem?}

In GA, one candidate solution is called a \textbf{chromosome}, and each entry is
a \textbf{gene}. Our chromosome is exactly the allocation list from Section~2:

\[
\text{chromosome} = [\,r_0, r_1, \dots, r_{399}\,],\qquad
\text{gene } r_i \in \{0,1,2,3,4\}.
\]

\begin{tcolorbox}[colback=blue!5!white,colframe=blue!50!black,title=GA is a natural fit for this problem]
Unlike PSO (which needs continuous numbers and a floor step), GA works
\textbf{directly on the discrete robot indices}. Gene $i$ literally \emph{is} the
robot assigned to flower $i$ --- no decoding needed. This is why an assignment
problem is a textbook use-case for GA.
\end{tcolorbox}

\subsection{Step 1 --- the population (and a smart head start)}
We start with a \textbf{population of 40 chromosomes}. Most are random, but we
give evolution a strong head start: a fraction (30\%) of the initial population
is \textbf{seeded} with a greedy ``nearest-depot'' allocation --- each flower
assigned to its \emph{closest} robot --- and then partly randomised for
diversity.

\begin{whybox}
\textbf{Heuristic seeding} is a standard, legitimate GA technique. It hands the
GA good \emph{spatial structure} to evolve from (short routes), while the fitness
function's balance term pushes the population toward \emph{balanced} loads. This
is how our GA fixes the exact weakness PSO had in 400 dimensions.
\end{whybox}

\subsection{Step 2 --- fitness (how good is each chromosome?)}
Every chromosome is scored with the \emph{same} fitness
function~\eqref{eq:fitness} used by PSO. Lower is better.

\subsection{Step 3 --- selection (choosing parents)}
We use \textbf{tournament selection}: pick a few chromosomes at random (we use
\textbf{3}) and let the fittest of that small group become a parent. Repeat to
get the second parent.

\begin{intuition}
Hold a mini-contest between 3 random candidates; the winner gets to reproduce.
Better chromosomes win more often, but even a weaker one occasionally sneaks
through --- which keeps healthy diversity. (This matches the ``roulette /
tournament / rank'' selection ideas from the lecture slides.)
\end{intuition}

\subsection{Step 4 --- crossover (combining two parents)}
We use \textbf{uniform crossover}: for each of the 400 genes, the child copies
that gene from parent~1 or parent~2 with equal probability.
\[
\text{child}_i =
\begin{cases}
\text{parent1}_i & \text{with probability } 0.5,\\
\text{parent2}_i & \text{with probability } 0.5.
\end{cases}
\]

\begin{intuition}
The child inherits some flower-assignments from mum and some from dad. If both
parents were good in \emph{different} regions of the greenhouse, the child can
combine the best of both --- this is the \textbf{exploitation} step that refines
good solutions. We apply crossover with probability $0.9$.
\end{intuition}

\subsection{Step 5 --- mutation (keeping diversity)}
We use \textbf{random-reset mutation}: each gene, with a small probability
(\textbf{2\%}), is reassigned to a random robot.

\begin{intuition}
A tiny random shake-up. Mutation is the \textbf{exploration} step: it introduces
brand-new gene values so the population never gets stuck copying the same few
solutions (``premature convergence,'' warned about in the slides). Too much
mutation = random search; too little = stagnation. $2\%$ is the sweet spot we
tuned.
\end{intuition}

\subsection{Step 6 --- elitism (never lose the best)}
The best \textbf{2} chromosomes are copied \emph{unchanged} into the next
generation.

\begin{whybox}
\textbf{Elitism} guarantees the best-so-far solution can never be destroyed by an
unlucky crossover or mutation. It makes the best-fitness curve
\emph{monotonically improving} (it can only get better or stay flat), which is
exactly what our convergence graph shows.
\end{whybox}

\subsection{The full GA procedure (step by step)}
\begin{enumerate}[nosep]
  \item Build a population of 40 chromosomes (30\% heuristic-seeded, rest random).
  \item \textbf{Repeat} for 120 generations:
  \begin{enumerate}[nosep]
    \item Score every chromosome with the fitness function~\eqref{eq:fitness}.
    \item Copy the best 2 (elitism) into the new generation.
    \item Fill the rest: tournament-select 2 parents $\to$ crossover (prob 0.9)
          $\to$ mutate (2\% per gene) $\to$ add the child.
    \item The new generation replaces the old one.
  \end{enumerate}
  \item Return the best chromosome ever seen.
\end{enumerate}

\begin{codebox}
This is \texttt{run\_ga\_allocation(...)} in \texttt{GeneticAlgorithm.py}.
Hyperparameters: \texttt{POP\_SIZE=40, N\_GEN=120, TOURNAMENT\_K=3, CX\_RATE=0.9,
MUT\_RATE=0.02, ELITE=2, SEED\_FRAC=0.30}. It imports
\texttt{generate\_greenhouse} and \texttt{evaluate} from \texttt{PSO.py} so both
allocators are judged on an identical environment with an identical fitness
function. A fixed seed makes it fully reproducible.
\end{codebox}

%=====================================================================
\section{PSO vs GA --- Side by Side}
%=====================================================================

\begin{table}[h]
\centering
\renewcommand{\arraystretch}{1.3}
\begin{tabular}{@{}p{3.4cm} p{5.1cm} p{5.1cm}@{}}
\toprule
& \textbf{PSO} & \textbf{GA} \\
\midrule
Nature's inspiration & Flock of birds & Biological evolution \\
A candidate is called & A particle & A chromosome \\
Encoding & 400 real numbers, floored to robot indices & 400 discrete robot indices (no decoding) \\
How solutions change & \emph{Move} through space via velocity (pbest + gbest) & \emph{Breed}: selection + crossover + mutation \\
Memory used & Personal best + global best position & The whole population + elitism \\
Exploration vs exploitation & Inertia + random $r_1,r_2$ & Mutation (explore) + crossover (exploit) \\
Key parameters & $w, c_1, c_2$, swarm size & Population, crossover \& mutation rate, elitism \\
Result (this project) & fitness $\approx 1.59$, dist $\approx 2450$ & fitness $\approx 1.43$, dist $\approx 2204$ (\textbf{better}) \\
Workload balance (std) & $\approx 2.1$ & $\approx 1.1$ (more balanced) \\
\bottomrule
\end{tabular}
\caption{PSO and GA solve the same allocation problem in different styles.}
\end{table}

\subsection{Why GA wins here}
GA has two structural advantages for this problem. First, its \textbf{discrete
encoding} matches the problem exactly (a flower goes to robot 2 \emph{or} 3,
never 2.5), avoiding the awkward continuous-to-discrete floor step. Second, our
\textbf{heuristic seeding} gives the GA a good starting point (nearest-depot
clusters), which crossover and selection then \emph{balance} under the fitness
function. Together these let GA reach a shorter total distance ($\approx 2204$ vs
$2450$) and a much more even workload (load std $\approx 1.1$ vs $2.1$), giving a
lower (better) overall fitness.

\begin{whybox}
\textbf{An honest note for your defense.} On the allocation \emph{map}, the GA's
colours still look intermixed --- much like PSO's. This is expected: the $400$
flowers form a few dense clusters, and keeping every robot's load near $80$
\emph{forces} those clusters to be shared between robots. The GA's advantage is
therefore \textbf{quantitative} (shorter total routes, tighter load balance,
lower fitness on the identical objective), \emph{not} a visually neater map.
Claiming otherwise would not survive a close look at the figure.
\end{whybox}

\begin{tcolorbox}[colback=green!5!white,colframe=green!45!black,title=Take-away]
Both are \textbf{population-based optimizers}: they keep many candidate
allocations and improve them over iterations. PSO \emph{moves} solutions using
personal/global-best memory; GA \emph{breeds} solutions using selection,
crossover and mutation. For our high-dimensional, spatial, discrete allocation
problem, GA's natural discrete encoding plus a heuristic head start is the better
fit --- and having \emph{both} lets us demonstrate a fair, meaningful comparison
in the report.
\end{tcolorbox}

%=====================================================================
\section{Where This Fits in the Whole Project}
%=====================================================================

\begin{center}
\begin{tabular}{@{}ll@{}}
\toprule
\textbf{Phase} & \textbf{Job} \\
\midrule
Phase 1 (this document) & PSO / GA decide \emph{which robot} gets \emph{which flowers}. \\
Phase 2 & Value Iteration \& Q-Learning navigate a robot cell-by-cell \\
         & through the grid to actually reach its flowers. \\
\bottomrule
\end{tabular}
\end{center}

\noindent
PSO and GA answer the \textbf{``who does what''} question. The Phase-2
reinforcement-learning algorithms then answer the \textbf{``how do I get
there''} question, using the assignments produced here. Each algorithm stays in
its own lane --- allocation and navigation are never mixed.

\end{document}
